\section{Learning Objectives}

After completing this chapter, readers will be able to:
\begin{itemize}
    \item Set up and use MLflow for experiment tracking
    \item Track model parameters, metrics, and artifacts
    \item Compare experiments and analyze results
    \item Manage model lifecycle with model registry
    \item Implement experiment tracking best practices
\end{itemize}

\section{Introduction to Experiment Tracking}

\subsection{Why Track Experiments?}

Systematic experiment tracking enables reproducibility, collaboration, and continuous improvement.

\section{MLflow Tutorial}

\begin{lstlisting}[style=python]
import mlflow
import mlflow.sklearn
from sklearn.ensemble import RandomForestClassifier

# Set tracking URI
mlflow.set_tracking_uri("http://localhost:5000")
mlflow.set_experiment("classification_experiments")

# Start MLflow run
with mlflow.start_run():
    # Log parameters
    mlflow.log_param("n_estimators", 100)
    mlflow.log_param("max_depth", 10)

    # Train model
    model = RandomForestClassifier(n_estimators=100, max_depth=10)
    model.fit(X_train, y_train)

    # Log metrics
    accuracy = model.score(X_test, y_test)
    mlflow.log_metric("accuracy", accuracy)

    # Log model
    mlflow.sklearn.log_model(model, "model")
\end{lstlisting}

\section{Model Registry}

% Placeholder content

\section{Chapter Summary}

% Placeholder summary

\section{Exercises}

\begin{exercise}
Set up MLflow tracking for a multi-model comparison experiment.
\end{exercise}
